\documentclass[11pt]{article}
\usepackage[margin=0.85in]{geometry}
\usepackage[T1]{fontenc}
\usepackage{lmodern,amsmath,graphicx,booktabs}
\pagestyle{empty}
\setlength{\parindent}{0pt}
\setlength{\parskip}{5pt}
\begin{document}
{\large\bfseries ACTSC 972: individual hedging experiment}\hfill Name: \rule{3.5cm}{0.4pt}

\textbf{Question and prediction.}
State the question and what you expected before inspecting test results.

\textbf{Design.}
State what changed, what stayed fixed, the training seed, test sample size,
which strategies were retrained, and the computational budget.

\textbf{Results.}
Replace this paragraph with one compact figure and/or table with units and
clear labels. The separate notebook PDF contains the full output and implementation checks.
For architecture comparisons, give parameter counts and training times.

\textbf{Interpretation and limitation.}
Explain the main finding and one limitation. Avoid interpreting a single
training seed as a general ranking of methods.

\textbf{Sources and reproducibility.}
Briefly identify external sources and refer to the notebook's
configuration and check results. Keep the complete report, including the
figure/table, within one page. You may use another editor and export a PDF.
Submit the complete notebook PDF separately; its length is not limited to one page.
\end{document}
